\documentclass[twoside, 10pt]{article}
\usepackage{geometry}
\geometry{bottom=4em,top=6em, inner=2.2cm, outer=4em,  headheight=\paperheight}
\usepackage[export]{adjustbox}
\usepackage{array}
\usepackage{amsmath}
\usepackage{amsfonts}
\usepackage{fancyhdr}
\usepackage{luacode}
\pagestyle{fancy}
\fancyhf{}
\lhead{Precalculus - BASE}
\chead{Properties of Exponents and Logarithms}
\directlua{pageNum = 1}
\rhead{Practice, Page \directlua{tex.print(pageNum); pageNum = pageNum + 1; if pageNum > 3 then pageNum = 1 end}}
\usepackage{lastpage}
\usepackage{xcolor}
\usepackage{enumitem}
\usepackage{pifont}
\usepackage{graphicx}
\graphicspath{{../img}}
\usepackage{pgfplots}
\pgfplotsset{compat=1.18}
\usepackage{tabularx}
\usepackage{polynom}

\newcommand{\R}{\mathbb R}
\newcommand{\e}{{\rm e}}
\newcommand{\pobr}[1]{\left\langle#1\right\rangle}
\newcommand{\norm}[1]{\lVert #1 \rVert}
\newcommand{\abs}[1]{\lvert #1 \rvert}

\DeclareMathOperator{\xd}{d\!}
\DeclareMathOperator{\proj}{proj}

\title{}
\date{}

\newcommand{\template}{
\directlua{tex.print(utils.genVersion())}
\noindent
{\large
First Name \rule{6em}{.1pt}\hspace{\stretch{1}} Last Name \rule{6em}{.1pt}\hspace{\stretch{1}} Date \rule{1.5em}{.1pt} -- \rule{1.5em}{.1pt} -- \rule{1.5em}{.1pt}\hspace{\stretch{1}} Period \rule{2em}{.1pt}\hspace{\stretch{1}} Score \rule{2em}{.1pt}
}
\vspace{1em}
}

\begin{luacode*}
utils = require("../utils")
\end{luacode*}

\begin{document}
\template

\noindent\textbf{Learning Target: Argue.}
\begin{itemize}
\item
I can prove the basic properties of exponents.
\item 
I can apply the appropriate properties of exponents to simplify an expression. 
\end{itemize}

\noindent{\bf Diagnostic Warm-up.}
Find the value of each of the following expressions.
\begin{enumerate}[itemsep=2em]
\item
$2^3 = \rule{5em}{0.4pt}$
\item 
$8^{2/3} = \rule{5em}{0.4pt}$
\item 
$2^{-1} = \rule{5em}{0.4pt}$
\item
$(2^{1/2})^4 = \rule{5em}{0.4pt}$
\item 
$2^{4/3}\cdot2^{2/3} = \rule{5em}{0.4pt}$
\item
$(4^3\cdot 3^6)^{1/6} = \rule{5em}{0.4pt}$
\end{enumerate}
\vspace{2em}


Let $n$ be a natural number, namely, a positive integer.
An expression of the form
\[
a^n
\]
means that the base $a$ is multiplied by itself $n$ times. That is,
\[
a^n = \underbrace{a \cdot a \cdot a \cdots a}_{n \text{ times}}.
\]
\begin{enumerate}[label=\arabic*.]
\item Prove $a^m \cdot a^n = a^{m + n}$.
\vspace{\stretch{1}}
\item Prove $(a^m)^n  = a^{m\cdot n}$
\vspace{\stretch{1}}
\clearpage
\item {\bf Exit ticket (Argue).}  Score \rule{4em}{0.4pt}

Exponents also have the following useful properties: Assume $m > n > 0$ and both are integers, then
\begin{enumerate}
\item
 $\displaystyle \frac{a^m}{a^n} = a^{m - n}$
 \item
 $(ab)^n = a^nb^n$
 \item
 $\displaystyle \left(\frac{a}{b}\right)^n =\frac{a^n}{b^n} $
\end{enumerate}
    Pick one to prove.
\clearpage

\item
Explain why $a^{-n}$ must be defined as $\displaystyle \frac{1}{a^n}$ in order to preserve the compatibility with the previously proved properties.
\vspace{\stretch{1}}

\item 
Explain why $a^{1/n}$ must be defined as $\displaystyle \sqrt[n]{a}$ in order to preserve the compatibility with the previously proved properties and the continuity of the exponential function. 
\vspace{\stretch{1}}
\clearpage

\item {\bf Exit Ticket (Argue).} Score \rule{4em}{.4pt}

Explain why $a^0 = 1$. You can use any properties that are proved previously.
\vspace{\stretch{1}}

\item {\bf Exit Ticket (Be Precise).} Score \rule{4em}{.4pt}

Find the value of each of the following expressions.
\begin{enumerate}[itemsep=2em]
\item
$3^3 = \rule{5em}{0.4pt}$
\item 
$9^{1/2} = \rule{5em}{0.4pt}$
\item 
$8^{-1/3} = \rule{5em}{0.4pt}$
\item
$(3^{1/2})^4 = \rule{5em}{0.4pt}$
\item 
$4^{1/4}\cdot4^{1/4} = \rule{5em}{0.4pt}$
\item
$\displaystyle\left(\frac{9^{1/4}}{4^{1/4}}\right)^2 = \rule{5em}{0.4pt}$
\end{enumerate}
\vspace{\stretch{1}}
\clearpage

\item
Complete the steps of a proof for the logarithm quotient rule: $\displaystyle \log_b\left(\frac{U}{V}\right) = \log_b{U} - \log_b{V}$. Start with two equations:
\[
b^x = U\hspace{8em}b^y = V
\]
\begin{enumerate}
\item
Rewrite both of these equations as logarithms.
\begin{align*}
\log_{\rule{2em}{.4pt}}\left(\rule{3em}{0.4pt}\right) &= \rule{4em}{.4pt}\\[1em]
\log_{\rule{2em}{.4pt}}\left(\rule{3em}{0.4pt}\right) &= \rule{4em}{.4pt}
\end{align*}
\item Divide the left sides of the original equations, and set the quotient equal to the quotient of
the right sides of the original equations.
\[
\rule{6em}{0.4pt} = \frac{U}{V}
\]
\item Combine the exponents on the left side of the equation so that it is written with a
single base.
\[
\rule{6em}{0.4pt} = \frac{U}{V}
\]
\item Rewrite the last equation as a logarithm.
\[
\log_{\rule{2em}{.4pt}}\left(\rule{3em}{0.4pt}\right) = \rule{4em}{.4pt}
\]
\item Combine the result in (i) and (iv) to derive the quotient rule of logarithms:\\[1em]
\begin{center}
\rule{20em}{.4pt}
\end{center}
\end{enumerate}
\item 
Prove the following {\bf Change of Base} formula: for any positive number $c$ (except 1),
\[
\log_a(b) = \frac{\log_c(b)}{\log_c(a)}
\]
\begin{enumerate}
\item
Start with the expression $\displaystyle \log_c(a^{\log_a(b)})$. Rewrite this expression in two ways.\\[1em]
\begin{enumerate}
\item
Use the power rule for logarithms. \rule{20em}{0.4pt}\\[1em]
\item Rewrite the expression $\displaystyle a^{\log_a(b)}$ so that it does not use a power.  \rule{15em}{0.4pt}
\end{enumerate}
\item Set the two expressions you wrote equal to each other. How can you turn this equation
into the change of base rule?\\[1em]

\hrulefill\\[1em]

\hrulefill\\[1em]

\hrulefill\\[1em]

\hrulefill\\[1em]

\hrulefill\\[1em]
\clearpage
\end{enumerate}
\item {\bf Exit Ticket (Argue).} Score \rule{4em}{.4pt}

Pick one to prove:

\begin{enumerate}
\item
$\log_b(U\cdot V) = \log_bU + \log_bV$
\item 
$\log_b(U^c) = c\log_b(U)$
\end{enumerate}
\end{enumerate}


\end{document}
\end{document}